\scanpage{99}
$\mathbb P(V_{P_n})(\mathbb C)$ satisfait \deleted{ainsi \unsure{que}} la propriété caractéristique de $\Sigma_{P_n}$, comme enveloppe-droite-projective de \deleted{$P_n$} $S_{P_n}$, ``\unsure{totale}'' compatible avec l'action de $D_{P_n}$.

Si par contre $n$ est pair, l'homomorphisme \deleted{\unsure{de}} $D_{\widetilde P_{2n}}\longrightarrow D_{P_n}$ ne scinde pas, car l'homomorphisme
\diagram{\begin{tikzcd}[ampersand replacement=\&,column sep=large]
D_{\widetilde P_{2n}}^+\arrow[r]\arrow[d,phantom,"\rotatebox{90}{$\simeq$}"]\&
D_{P_n}^+\arrow[d,phantom,"\rotatebox{90}{$\simeq$}"]\\
\mathbb Z/2n\&\mathbb Z/n
\end{tikzcd}}
ne scinde pas (il est isom. à l'hom. de passage au quotient canonique $\mathbb Z/2n\mathbin{-}\mathbb Z/n$, qui en passant au quotient par $\deleted{\unsure{(2m).\mathbb Z}}\ m.\mathbb Z$ (où $n=2m$) donne $\mathbb Z/4\mathbb Z\longrightarrow\mathbb Z/2\mathbb Z$, qui n'est pas scindée. \deleted{Donc} Comme l'extension (8) de $D_{P_n}$ par $\{\pm1\}$ n'est autre que l'extension réciproque de
\diagram{\begin{tikzcd}[ampersand replacement=\&,column sep=small]
1\arrow[r,no head]\&\{\pm1\}\arrow[r]\&
\deleted{\unsure{\mathrm{Gl}^{?}}}\ \unsure{\mathrm{Sl}|\mathrm{Gl}}(V_{\widetilde P_{2n}})
\arrow[r]\&
\underset{\substack{\Vert\\\operatorname{Aut}(\Sigma_{P_n})}}{\mathrm{GP}(V_{\widetilde P_{2n}})}
\arrow[r,no head]\&1
\end{tikzcd}}
par $D_{P_n}\longrightarrow\operatorname{Aut}(\Sigma_{P_n})$, cela montre que l'opération de $D_{P_n}$ sur $\Sigma_{P_n}$ ne provient pas d'une représentation vectorielle — \deleted{du moins} ce qui revient au $\unsure{=|\overline m}$, qu'elle ne se remonte pas

\margin{98}
